%======================================================================
\section{The grid host}\label{sec:grid}
%======================================================================

This appendix is a sketch: its conclusions are conditional on proofs that we only outline, and they are stated in the indicative with that
understanding. Nothing elsewhere in the paper depends on them. Placing the lamps on configurations of the cells of a grid, rather than of the
points of $Y$, gives a stronger lower bound with the same upper bound. It also gives a nearly linear exponent $r^{1-o(1)}$ from Lee's
exponential Dehn bound alone.

\begin{remark}[The grid host]\label{rem:grid}
Let $I=Y\times Y$, let $T=H\times H$ act on $I$ coordinatewise, and let
\[
\Gamma_2=\Bigl(\bigoplus_{v\in V_2}S_3\Bigr)\rtimes\bigl(V_2\rtimes(\GL_{\rm fin}(I)\rtimes T)\bigr),\qquad V_2=\F_2^{(I)} .
\]
It is locally finite-by-$\Z^4$ and generated by seven elements: the generators $p_1,q_1,p_2,q_2$ of the two copies of $H$, the translation $t$
by the cell $o=((1,1),(1,1))$, the element $c=d\,ab$, where $d$ is the transvection that adds $e_o$ to $e_{((1,2),(1,2))}$, and $a$.
Transvections within a row or a column are not needed as generators, since commutators $[E_{x\leftarrow y},E_{y\leftarrow z}]=E_{x\leftarrow z}$
of transvections of the third kind produce them. A certificate of $77$ relations, of total length $5445$ and maximum length $224$, is listed in
\texttt{verification/relators.txt} and checked by \texttt{verification/relators\_check.py}. It has the families of Table~\ref{tab:certificate},
with two copies of the Houghton relators and their commutation, and with stabilizer and action relations for three kinds of transvection: within
a row, within a column, and neither.

Transport now happens at scale $\sqrt N$. $N$ cells fit in a square of side $\lceil\sqrt N\,\rceil$, so the transporters have length $O(\sqrt N)$,
and the local fillings cost $C(\delta_H(C\sqrt N)+N)$; the term $N$ counts the swaps of letters from the two copies of $H$ in a change of
transporter. The clearing argument of Theorem~\ref{thm:pair-area} then bounds the pair areas by $CN^2(\delta_H(C\sqrt N)+N)$, which is at most
$CN^{2+p_H/2}$ by Proposition~\ref{prop:dehn}(b) and at most $e^{C\sqrt N}$ by Proposition~\ref{prop:dehn}(a). As in the proof of
Theorem~\ref{thm:main-group}, the chunk $E_2$ of this certificate therefore has
\[
\log\sigma_{E_2}(r)\ \ge\ c\,\frac{r}{(\log r)^{2+p_H/2}},\qquad 2+\tfrac{p_H}2<7.38,
\]
and, using only Lee's bound, $\log\sigma_{E_2}(r)\ge c\,r\,e^{-C\sqrt{\log r}}$; unitary models need $\log D\ge c_\Delta e^{-2}/(\log(1/e))^{4+p_H}$.
In the other direction, hashing the configurations directly, as in Proposition~\ref{prop:upper}, needs a configuration space of dimension of
order $r^2$. Hashing the cells first avoids this. A cell $(y_1,y_2)$ gets the label $a_1(f_1^{-1}y_1)+a_2(f_2^{-1}y_2)\in\F_q$. Here the $f_j$ are
positions in finite F{\o}lner charts of $H$ of size $e^{O(r\log r)}$, inverses of right F{\o}lner sets because the charts move by left
multiplication, and the $a_j$ are random tables on $O(r)$ points. The label is covariant under the charts, and it separates the finitely many
cells of a chunk with probability $1-O(1/q)$. The configurations are then compressed linearly to $\F_2^{\F_q}$ and hashed affinely as in
Proposition~\ref{prop:upper}. With $q$ of order $\ell^2r$, for a chunk of words of length at most $\ell$, and the other parameters of order $r$,
this gives $\log\sigma_{E'}(r)\le C_{E'}r\log(2+r)$ for every chunk $E'$ of $\Gamma_2$. The same holds for the $d$-dimensional grid host, built
in the same way on the cells of $Y^d$ with $H^d$ acting, for every fixed $d$. Hence $\sigma_{E_2}(r)=\exp(r^{1+o(1)})$ as well, and no
fixed-dimensional grid host exceeds $\exp(r^{1+o(1)})$. The proofs follow those of Sections~\ref{sec:fillings}--\ref{sec:profile} with the
changes indicated.
\end{remark}

\begin{remark}[The gadget channel of the grid host]\label{rem:grid-channel}
Adding $b$ and $j$ and the four words of Section~\ref{sec:lamp-channel} to the certificate of $\Gamma_2$ gives $s=9$, $r_0=81$ and $\Lambda=5462$,
so the gadget channel $\Phi_2$ of $\Gamma_2$ acts on $M_d$ with
\[
d=1+4\cdot9+3\cdot5462-4\cdot81=16099,
\]
and has $m=224$. Of its $81$ words, $78$ have length at least three, so by the count of Section~\ref{sec:lamp-channel} its Choi rank is at most
$1+18+5300-78-3=5238$. With the transport of Remark~\ref{rem:grid}, words of length $O(N^{3/2})$ and pair areas
$O(N^2(\delta_H(C\sqrt N)+N))$, the hypothesis of Corollary~\ref{cor:amplification}(a) holds with the exponent $2+p_H/2$, and part~(b) of that
corollary holds with Lee's exponential Dehn bound~\cite{Lee2012} alone. The proof of Theorem~\ref{thm:main-channel} then gives three bounds for
$\Phi_2$. At purity $O(\log n)$ the memory is at least $c\,n/(\log n)^{p_H+5}$, where $p_H+5<15.75$, and at least $n\,2^{-C\sqrt{\log n}}$ using
only Lee's bound. At the rank rate it is at least $c\sqrt n/(\log n)^{2+p_H/2}$, whatever the purity. By the two-stage hashing of
Remark~\ref{rem:grid}, $\log\sigma(r)=O(r\log r)$ for $\Gamma_2$ as well, which gives an upper bound of order $n\log n$. For $\Gamma$ itself,
Lee's bound alone gives pair areas $2^{O(N)}$, and the argument gives only memory $n^{c}$ for some $c>0$.
\end{remark}
